\documentclass[12pt,a4paper]{article}

\usepackage[utf8]{inputenc}
\usepackage[T1]{fontenc}
\usepackage[brazil]{babel}
\usepackage[margin=2.5cm]{geometry}
\usepackage{array}
\usepackage{longtable}
\usepackage{booktabs}
\usepackage{xcolor}
\usepackage{hyperref}
\usepackage{titlesec}
\usepackage{enumitem}
\usepackage{tocloft}
\usepackage{graphicx}
\usepackage{float}

\hypersetup{
    colorlinks=true,
    linkcolor=black,
    urlcolor=blue,
    citecolor=black,
    pdftitle={RepMax - Analise e Defesa de Projeto},
    pdfauthor={Artur Pereira e Zadoque Carneiro}
}

\titleformat{\section}{\normalfont\Large\bfseries}{\thesection}{1em}{}
\titleformat{\subsection}{\normalfont\large\bfseries}{\thesubsection}{1em}{}

\newcolumntype{L}[1]{>{\raggedright\arraybackslash}p{#1}}

\title{\textbf{RepMax}\\[0.2cm]
\Large Sistema de Log Book de Treino de Academia\\[0.2cm]
\large Documentação de Análise, Requisitos e Arquitetura do Sistema}
\author{Artur Pereira e Zadoque Carneiro}
\date{Disciplina: POODev}

\begin{document}

\maketitle
\tableofcontents
\newpage

\section{Introdução}

O crescimento contínuo da prática de atividades físicas no Brasil, evidenciado pelos
dados da Sociedade Brasileira de Endocrinologia e Metabologia (SBEM)\footnote{SOCIEDADE BRASILEIRA DE ENDOCRINOLOGIA E METABOLOGIA (SBEM). \textit{Dados sobre a prática de atividade física no Brasil}. [link oficial e ano da pesquisa --- campo editável, a ser preenchido com a fonte efetivamente consultada].}, impulsiona a
demanda por soluções digitais capazes de organizar e acompanhar a evolução de treinos de
forma estruturada. Assim, o \textbf{RepMax} surge como uma plataforma web voltada ao registro
de treinos de musculação, à visualização de progresso e à integração entre alunos e
personal trainers, somando mais tecnologia ao meio fitness.


\subsection{Problema de Negócio}

Observam-se, entre praticantes de musculação, as seguintes dificuldades recorrentes:

\begin{itemize}
    \item Ausência de um histórico organizado de séries, repetições e cargas utilizadas em cada treino;
    \item Dificuldade de identificar se um exercício está em progressão, estagnação ou regressão de desempenho;
    \item Falta de suporte para técnicas de intensificação de treino (bi-set, tri-set, monster set e drop set), que hoje são anotadas de forma manual e pouco estruturada;
    \item Dúvidas frequentes sobre a execução correta de determinados exercícios, sem uma referência disponível no momento do treino;
    \item Ausência de um canal integrado entre aluno e personal trainer, que hoje depende de aplicativos de mensagens externos (como WhatsApp) e planilhas avulsas.
\end{itemize}

\subsection{Proposta do Trabalho}
Embora esse projeto tenha como objetivo principal resolver o problema de negócio e agragar valor para o cliente, há também o propósito de aprendizado de conceitos. Por isso, esse trabalho propõe o desenvolvimento de uma API RESTful para um sistema de log book de treino, denominado \textbf{RepMax}. E uma interface WEB usando: React + Typescript + Tailwind CSS + React Router (\textbr{Data Mode}) + Vite + fetch() para interagir com ela (Mais abaixo tem a justificativa para o uso de cada parte da interface). A solução tem como objetivo permitir que o usuário registre seus treinos de forma estruturada, acompanhe sua progressão por meio de indicadores, tais como: a repetição máxima estimada (1RM),  técnicas avançadas de intensidade e, em uma segunda versão do sistema, viabilize a atuação de um personal trainer na montagem de treinos e no acompanhamento de alunos.

\section{Descrição Geral do Sistema}

\subsection{API RESTful}

Ela será responsável por gerenciar o histórico de treino de usuários. De forma geral, o sistema deverá permitir:

\begin{itemize}
    \item Cadastrar e autenticar usuários (alunos, personais e admins);
    \item Cadastrar exercícios, com grupo muscular e instruções de execução;
    \item Registrar treinos, associando exercícios, séries, repetições e cargas;
    \item Registrar séries em formatos de intensificação: bi-set, tri-set, monster set e drop set;
    \item Calcular a repetição máxima estimada (1RM) a partir de carga e repetições;
    \item Consultar o histórico e a evolução de desempenho por exercício;
    \item (Versão 2) Vincular aluno e personal trainer, permitindo que o personal monte treinos, o aluno registre feedback e ambos se comuniquem dentro do próprio sistema.
\end{itemize}

\subsection{Interface WEB:}
O conjunto de tecnologia (stack) foi escolhido para viabilizar o aprendizado de conceitos como: rotas, cache, segurança, server-side e client-side.
Vite -> Ferramenta responsável pelo ambiente dev e processo de build do front-end
React -> Biblioteca de UI, para construção e navegação na interface.
Tailwind CSS -> Para estilização.
React Router no modo Data -> Não faz mágica sozinho, mas traz ferramentas para eu aprender e usar como Rotas, Loader, Action e Estados de navegação.
fetch() -> Para fazer requisição à API RESTful.

Esse conjunto permite que tenhamos um equilíbrio entre conceitos a aprender e aplicar e prazo de entrega. 
\end{document}
